\chapter{2009年新课标卷（文）}

\section{单选题}

\begin{problem}
已知集合 \(A = \{1, 3, 5, 7, 9\}\)，\(B = \{0, 3, 6, 9, 12\}\)，则 \(A \cap B =\)（\quad）

\choices
    {\( \{3, 5\} \)}
    {\(\{3,6\}\)}
    {\(\{3,7\}\)}
    {\(\{3,9\}\)}
\end{problem}

\begin{answer}
D.
\end{answer}

\begin{solution}
交集只包含同时属于两个集合的元素，因此
\[
A\cap B=\{3,9\}
\]
故选 D.
\end{solution}

\begin{problem}
复数 \(\frac{3 + 2\i}{2 - 3\i} =\)（\quad）

\choices
    {\(1\)}
    {\(-1\)}
    {\(\i\)}
    {\(-\i\)}
\end{problem}

\begin{answer}
C.
\end{answer}

\begin{solution}
分子分母同乘以分母的共轭复数，得
\[
\frac{3+2\i}{2-3\i}
=\frac{(3+2\i)(2+3\i)}{(2-3\i)(2+3\i)}
=\frac{6+9\i+4\i+6\i^{2}}{13}
=\frac{13\i}{13}=\i
\]
故选 C.
\end{solution}

\begin{problem}
对变量 \(x\)，\(y\) 有观测数据 \((x_{i}, y_{i})\)（\(i = 1\)，\(2\)，\(\dots\)，\(10\)），得散点图 1；对变量 \(u\)，\(v\) 有观测数据 \((u_{i}, v_{i})\)（\(i = 1\)，\(2\)，\(\dots\)，\(10\)），得散点图 2. 由这两个散点图可以判断（\quad）

\bitmapfigure[float=inline,max width=0.68\linewidth]{img/2009/new_curriculum_liberal/q3_fig1.png}

\choices
    {变量 \(x\) 与 \(y\) 正相关，\(u\) 与 \(v\) 正相关}
    {变量 \(x\) 与 \(y\) 正相关，\(u\) 与 \(v\) 负相关}
    {变量 \(x\) 与 \(y\) 负相关，\(u\) 与 \(v\) 正相关}
    {变量 \(x\) 与 \(y\) 负相关，\(u\) 与 \(v\) 负相关}
\end{problem}

\begin{answer}
C.
\end{answer}

\begin{solution}
由散点图 1 可见，随着 \(x\) 增大，\(y\) 总体减小，所以 \(x\) 与 \(y\) 负相关；由散点图 2 可见，随着 \(u\) 增大，\(v\) 总体增大，所以 \(u\) 与 \(v\) 正相关.

因此选 C.
\end{solution}

\begin{problem}
有四个关于三角函数的命题：

\begin{circlelist}
\item \(p_{1}\)：\(\exists x \in \R\)，\(\sin^{2}\frac{x}{2} + \cos^{2}\frac{x}{2} = \frac{1}{2}\)；
\item \(p_{2}\)：\(\exists x\)，\(y \in \R\)，\(\sin(x - y) = \sin x - \sin y\)；
\item \(p_{3}\)：\(\forall x \in [0,\pi]\)，\(\sqrt{\frac{1 - \cos 2x}{2}} = \sin x\)；
\item \(p_{4}\)：\(\sin x = \cos y \Rightarrow x + y = \frac{\pi}{2}\).
\end{circlelist}

其中假命题的是（\quad）

\choices
    {\(p_1\)，\(p_4\)}
    {\(p_{2}\)，\(p_{4}\)}
    {\(p_1\)，\(p_3\)}
    {\(p_2\)，\(p_3\)}
\end{problem}

\begin{answer}
A.
\end{answer}

\begin{solution}
对于 \(p_{1}\)，由同角三角函数基本关系式，
\[
\sin^{2}\frac{x}{2}+\cos^{2}\frac{x}{2}=1
\]
不可能等于 \(\frac12\)，所以 \(p_{1}\) 为假命题.

对于 \(p_{2}\)，取 \(y=0\)，则
\[
\sin(x-y)=\sin x=\sin x-\sin0
\]
所以存在满足条件的 \(x\)，\(y\)，\(p_{2}\) 为真命题.

对于 \(x\in[0,\pi]\)，有 \(\sin x\geqslant0\)，且
\[
\sqrt{\frac{1-\cos2x}{2}}
=\sqrt{\sin^{2}x}=\sin x
\]
所以 \(p_{3}\) 为真命题.

命题 \(p_{4}\) 不成立. 例如取 \(x=\frac{\pi}{6}\)，\(y=\frac{\pi}{3}+2\pi\)，则 \(\sin x=\cos y\)，但
\[
x+y=\frac{5\pi}{2}\ne\frac{\pi}{2}
\]
因此假命题是 \(p_{1}\)，\(p_{4}\)，故选 A.
\end{solution}

\begin{problem}
已知圆 \(C_1: (x + 1)^2 + (y - 1)^2 = 1\)，圆 \(C_2\) 与圆 \(C_1\) 关于直线 \(x - y - 1 = 0\) 对称，则圆 \(C_2\) 的方程为（\quad）

\choices
    {\((x + 2)^{2} + (y - 2)^{2} = 1\)}
    {\((x - 2)^{2} + (y + 2)^{2} = 1\)}
    {\((x + 2)^{2} + (y + 2)^{2} = 1\)}
    {\((x - 2)^{2} + (y - 2)^{2} = 1\)}
\end{problem}

\begin{answer}
B.
\end{answer}

\begin{solution}
圆 \(C_{1}\) 的圆心为 \((-1,1)\)，半径为 \(1\). 点 \((-1,1)\) 关于直线 \(x-y-1=0\) 的对称点为 \((2,-2)\)：事实上，过 \((-1,1)\) 作该直线的垂线，其垂足为 \(\left(\frac12,-\frac12\right)\)，故对称点为
\[
2\left(\frac12,-\frac12\right)-(-1,1)=(2,-2)
\]
对称变换保持半径不变，所以
\[
C_{2}:(x-2)^{2}+(y+2)^{2}=1
\]
故选 B.
\end{solution}

\begin{problem}
设 \(x\)，\(y\) 满足 \(\begin{cases}
2x + y \geqslant 4,\\
x - y \geqslant -1,\\
x - 2y \leqslant 2,
\end{cases}\) 则 \(z = x + y\)（\quad）

\choices
    {有最小值 \(2\)，最大值 \(3\)}
    {有最小值 \(2\)，无最大值}
    {有最大值 \(3\)，无最小值}
    {既无最小值，也无最大值}
\end{problem}

\begin{answer}
B.
\end{answer}

\begin{solution}
由前两个约束条件和第三个约束条件分别构造
\[
z-2=\frac35(2x+y-4)+\frac15(2-x+2y)
\]
其中 \(2x+y-4\geqslant0\)，\(2-x+2y=2-(x-2y)\geqslant0\)，所以 \(z\geqslant2\).

当 \((x,y)=(2,0)\) 时三个约束均满足，且 \(z=2\)，故 \(z\) 的最小值为 \(2\).

另一方面，取 \(x=y=t\)，其中 \(t\geqslant\frac43\)，则三个约束均满足，而 \(z=2t\) 可任意增大，因此 \(z\) 无最大值.

故选 B.
\end{solution}

\begin{problem}
已知 \(\bs{a} = (-3, 2)\)，\(\bs{b} = (-1, 0)\)，向量 \(\lambda \bs{a} + \bs{b}\) 与 \(\bs{a} - 2\bs{b}\) 垂直，则实数 \(\lambda\) 的值为（\quad）

\choices
    {\( -\frac{1}{7} \)}
    {\( \frac{1}{7} \)}
    {\(-\frac{1}{6}\)}
    {\( \frac{1}{6} \)}
\end{problem}

\begin{answer}
A.
\end{answer}

\begin{solution}
\(\lambda\bs{a}+\bs{b}=(-3\lambda-1,2\lambda)\)，\(\bs{a}-2\bs{b}=(-1,2)\).
两向量垂直，当且仅当它们的数量积为零，因此
\[
(-3\lambda-1)(-1)+2\lambda\cdot2=0
\]
即 \(7\lambda+1=0\)，解得
\[
\lambda=-\frac17
\]
故选 A.
\end{solution}

\begin{problem}
等差数列 \(\{a_{n}\}\) 前 \(n\) 项和为 \(S_{n}\). 已知 \(a_{m-1} + a_{m+1} - a_{m}^{2} = 0\)，\(S_{2m-1} = 38\)，则 \(m =\)（\quad）

\choices
    {\(38\)}
    {\(20\)}
    {\(10\)}
    {\(9\)}
\end{problem}

\begin{answer}
C.
\end{answer}

\begin{solution}
等差数列满足
\[
a_{m-1}+a_{m+1}=2a_{m}
\]
代入已知条件得
\[
2a_{m}-a_{m}^{2}=0
\]
所以 \(a_{m}=0\) 或 \(a_{m}=2\). 又因为 \(S_{2m-1}\) 中有 \(2m-1\) 项，且关于中间项 \(a_m\) 对称，故
\[
S_{2m-1}=(2m-1)a_m=38
\]
于是 \(a_m\ne0\)，从而 \(a_m=2\)，并且
\(2(2m-1)=38\)，\(m=10\).
故选 C.
\end{solution}

\begin{problem}
如图，正方体 \( ABCD - A_{1}B_{1}C_{1}D_{1} \) 的棱长为 \(1\)，线段 \(B_{1}D_{1}\) 上有两个动点 \(E\)，\(F\)，且 \( EF = \frac{\sqrt{2}}{2} \)，则下列结论中错误的是（\quad）

\bitmapfigure[float=inline,max width=0.36\linewidth]{img/2009/new_curriculum_liberal/q9_fig1.png}

\choices
    {\(AC \perp BE\)}
    {\(EF \parallel\) 平面 \(ABCD\)}
    {三棱锥 \(A - BEF\) 的体积为定值}
    {异面直线 \(AE\)，\(BF\) 所成的角为定值}
\end{problem}

\begin{answer}
D.
\end{answer}

\begin{solution}
建立单位正方体坐标系，取
\(D=(0,0,0)\)，\(A=(1,0,0)\)，\(C=(0,1,0)\)，\(B=(1,1,0)\)
顶面各点的 \(z\) 坐标均为 \(1\). 设
\(E=(t,t,1)\)，\(F=\left(t+\frac12,t+\frac12,1\right)\)，\(0\leqslant t\leqslant\frac12\)
这正表示 \(E\)，\(F\) 在 \(B_{1}D_{1}\) 上且 \(EF=\frac{\sqrt2}{2}\). 另一种先后次序只需交换 \(E\)，\(F\)，不影响下列前两个结论.

因为
\(\overrightarrow{AC}=(-1,1,0)\)，\(\overrightarrow{BE}=(t-1,t-1,1)\)
所以 \(\overrightarrow{AC}\cdot\overrightarrow{BE}=0\)，即 \(AC\perp BE\).

又 \(\overrightarrow{EF}=(\frac12,\frac12,0)\) 平行于底面 \(ABCD\)，所以 \(EF\parallel\) 平面 \(ABCD\).

直线 \(EF\) 在平面 \(z=1\) 内，点 \(B\) 到直线 \(EF\) 的距离为 \(1\). 因而
\[
S_{\triangle BEF}=\frac12\cdot\frac{\sqrt2}{2}\cdot1=\frac{\sqrt2}{4}
\]
平面 \(BEF\) 的方程为 \(x=y\)，点 \(A=(1,0,0)\) 到该平面的距离为
\[
\frac{|1-0|}{\sqrt{1^{2}+(-1)^{2}}}=\frac{1}{\sqrt2}
\]
所以
\[
V_{A-BEF}=\frac13\cdot\frac{\sqrt2}{4}\cdot\frac1{\sqrt2}=\frac1{12}
\]
为定值.

最后，取 \(\overrightarrow{AE}\) 和 \(\overrightarrow{BF}\) 为两异面直线的方向向量，则
\(\overrightarrow{AE}=(t-1,t,1)\)，\(\overrightarrow{BF}=\left(t-\frac12,t-\frac12,1\right)\).
令 \(s=2t^{2}-2t\)，则两向量夹角的余弦为
\[
\frac{\overrightarrow{AE}\cdot\overrightarrow{BF}}
{|\overrightarrow{AE}|\,|\overrightarrow{BF}|}
=\frac{s+\frac32}{\sqrt{(s+2)(s+\frac32)}}
\]
当 \(t=0\) 时该值为 \(\frac{\sqrt3}{2}\)，当 \(t=\frac12\) 时该值为 \(\sqrt{\frac23}\)，两者不相等，故异面直线 \(AE\)，\(BF\) 所成的角不为定值.

因此错误的结论是 D.
\end{solution}

\begin{problem}
如果执行如图的程序框图，输入 \(x = -2\)，\(h = 0.5\)，那么输出的各个数的和等于（\quad）

\bitmapfigure[float=inline,max width=0.72\linewidth]{img/2009/new_curriculum_liberal/q10_fig1.png}
\FloatBarrier

\choices
    {\(3\)}
    {\(3.5\)}
    {\(4\)}
    {\(4.5\)}
\end{problem}

\begin{answer}
B.
\end{answer}

\begin{solution}
按照程序框图依次输入 \(x=-2\)，\(-1.5\)，\(-1\)，\(-0.5\)，\(0\)，\(0.5\)，\(1\)，\(1.5\)，\(2\). 当 \(x<0\) 时输出 \(y=0\)；当 \(0\leqslant x<1\) 时输出 \(y=x\)；当 \(x\geqslant1\) 时输出 \(y=1\). 因此输出依次为
\(0\)，\(0\)，\(0\)，\(0\)，\(0\)，\(0.5\)，\(1\)，\(1\)，\(1\).
当 \(x=2\) 时满足结束条件 \(x\geqslant2\)，所以输出各数之和为
\[
0+0+0+0+0+0.5+1+1+1=3.5
\]
故选 B.
\end{solution}

\begin{problem}
一个棱锥的三视图如图，则该棱锥的全面积（单位：\(\mathrm{cm}^2\)）为（\quad）

\bitmapfigure[float=inline,max width=0.5\linewidth]{img/2009/new_curriculum_liberal/q11_fig1.png}
\FloatBarrier

\choices
    {\(48 + 12\sqrt{2}\)}
    {\(48 + 24\sqrt{2}\)}
    {\(36 + 12\sqrt{2}\)}
    {\(36 + 24\sqrt{2}\)}
\end{problem}

\begin{answer}
A.
\end{answer}

\begin{solution}
由三视图可知，棱锥的底面是直角等腰三角形，两直角边均为 \(6\)，且棱锥的高为 \(4\)，顶点在底面上的射影为斜边中点. 因而底面积为
\[
S_{\text{底}}=\frac12\cdot6\cdot6=18
\]

底面两条直角边对应的两个侧面，其底边均为 \(6\). 顶点射影到这两条直角边的距离均为 \(3\)，所以这两个侧面的高均为
\[
\sqrt{3^{2}+4^{2}}=5
\]
面积各为 \(\frac12\cdot6\cdot5=15\). 斜边长为 \(6\sqrt2\)，顶点射影在斜边上，因此对应侧面的高为棱锥的高 \(4\)，面积为
\[
\frac12\cdot6\sqrt2\cdot4=12\sqrt2
\]
所以棱锥的全面积为
\[
18+15+15+12\sqrt2=48+12\sqrt2
\]
故选 A.
\end{solution}

\begin{problem}
用 \(\min \{a, b, c\}\) 表示 \(a\)，\(b\)，\(c\) 三个数中的最小值. 设 \(f(x) = \min \{2^x, x + 2, 10 - x\}\)（\(x \geqslant 0\)），则 \(f(x)\) 的最大值为（\quad）

\choices
    {\(4\)}
    {\(5\)}
    {\(6\)}
    {\(7\)}
\end{problem}

\begin{answer}
C.
\end{answer}

\begin{solution}
因为
\[
f(x)\leqslant\min\{x+2,10-x\}
\]
而 \((x+2)+(10-x)=12\)，所以两个数的较小者不超过 \(6\)，从而 \(f(x)\leqslant6\).

取 \(x=4\)，有
\(2^{x}=16\)，\(x+2=6\)，\(10-x=6\)
故 \(f(4)=6\). 因此 \(f(x)\) 的最大值为 \(6\)，故选 C.
\end{solution}

\section{填空题}

\begin{problem}
曲线 \( y = x \e^{x} + 2x + 1 \) 在点 \((0,1)\) 处的切线方程为\fillinblank{}.
\end{problem}

\begin{answer}
\(y=3x+1\).
\end{answer}

\begin{solution}
令 \(f(x)=x\e^{x}+2x+1\)，则
\[
f'(x)=\e^{x}+x\e^{x}+2=(x+1)\e^{x}+2
\]
所以 \(f(0)=1\)，\(f'(0)=3\). 由点斜式，切线方程为
\[
y-1=3(x-0)
\]
即 \(y=3x+1\).
\end{solution}

\begin{problem}
已知抛物线 \(C\) 的顶点坐标为原点，焦点在 \(x\) 轴上，直线 \(y = x\) 与抛物线 \(C\) 交于 \(A\)，\(B\) 两点，若 \(P(2, 2)\) 为 \(AB\) 的中点，则抛物线 \(C\) 的方程为\fillinblank{}.
\end{problem}

\begin{answer}
\(y^{2}=4x\).
\end{answer}

\begin{solution}
抛物线顶点在原点且焦点在 \(x\) 轴上，可设其方程为
\(y^{2}=4ax\)，\(a\ne0\).
直线 \(y=x\) 与抛物线联立，得
\(x^{2}=4ax\)，\(x(x-4a)=0\).
所以交点为 \((0,0)\) 和 \((4a,4a)\)，其中点为 \((2a,2a)\). 由中点为 \(P(2,2)\)，得 \(2a=2\)，即 \(a=1\). 因而抛物线方程为
\[
y^{2}=4x
\]
\end{solution}

\begin{problem}
等比数列 \(a_{n}\) 的公比 \(q > 0\)，已知 \(a_{2} = 1\)，\(a_{n+2} + a_{n+1} = 6a_{n}\)，则 \(a_{n}\) 的前 \(4\) 项和 \(S_{4} =\)\fillinblank{}.
\end{problem}

\begin{answer}
\(\dfrac{15}{2}\).
\end{answer}

\begin{solution}
设等比数列公比为 \(q\)，则由 \(a_{2}=1\) 得
\(a_{1}=\frac1q\)，\(a_{3}=q\).
在递推关系中取 \(n=1\)，有
\[
a_{3}+a_{2}=6a_{1}
\]
即
\(q+1=\frac6q\)，\(q^{2}+q-6=0\).
因 \(q>0\)，所以 \(q=2\). 因此前四项为
\(\frac12\)，\(1\)，\(2\)，\(4\)
从而
\[
S_{4}=\frac12+1+2+4=\frac{15}{2}
\]
\end{solution}

\begin{problem}
已知函数 \( f(x)=2\sin(\omega x+\varphi) \) 的图象如图所示，则 \( f\left(\frac{7\pi}{12}\right)= \)\fillinblank{}.

\bitmapfigure[max width=0.65\linewidth]{img/2009/new_curriculum_liberal/q16_fig1.png}
\end{problem}

\begin{answer}
\(0\).
\end{answer}

\begin{solution}
从图象可见，振幅为 \(2\)，相邻两个正最大值的横坐标相差 \(\frac{2\pi}{3}\)，所以 \(\omega=3\). 图象在 \(x=\frac\pi4\) 处从下向上过 \(x\) 轴，故可取
\(3\cdot\frac\pi4+\varphi=0\)，\(\varphi=-\frac{3\pi}{4}\).
于是
\[
f\left(\frac{7\pi}{12}\right)
=2\sin\left(3\cdot\frac{7\pi}{12}-\frac{3\pi}{4}\right)
=2\sin\pi=0
\]
\end{solution}

\section{解答题}

\begin{problem}
如图，为了解某海域海底构造，在海平面内一条直线上的 \(A\)，\(B\)，\(C\) 三点进行测量，已知 \(AB = 50 \mathrm{~m}\)，\(BC = 120 \mathrm{~m}\)，于 \(A\) 处测得水深 \(AD = 80 \mathrm{~m}\)，于 \(B\) 处测得水深 \(BE = 200 \mathrm{~m}\)，于 \(C\) 处测得水深 \(CF = 110 \mathrm{~m}\)，求 \(\angle DEF\) 的余弦值.

\bitmapfigure[float=inline,max width=0.38\linewidth]{img/2009/new_curriculum_liberal/q17_fig1.png}
\FloatBarrier
\end{problem}

\begin{solution}
建立空间直角坐标系，使海平面为 \(z=0\) 平面，且 \(A\)，\(B\)，\(C\) 依次位于 \(x\) 轴上. 取竖直向上为 \(z\) 轴正方向，则
\(A=(0,0,0)\)，\(B=(50,0,0)\)，\(C=(170,0,0)\)
以及
\(D=(0,0,-80)\)，\(E=(50,0,-200)\)，\(F=(170,0,-110)\).
因此
\(\overrightarrow{ED}=(-50,0,120)\)，\(\overrightarrow{EF}=(120,0,90)\).
从而
\[
\overrightarrow{ED}\cdot\overrightarrow{EF}=(-50)\times120+120\times90=4800
\]
\(|\overrightarrow{ED}|=\sqrt{50^2+120^2}=130\)，\(|\overrightarrow{EF}|=\sqrt{120^2+90^2}=150\).
所以
\[
\cos\angle DEF=\frac{\overrightarrow{ED}\cdot\overrightarrow{EF}}{|\overrightarrow{ED}|\,|\overrightarrow{EF}|}=\frac{4800}{130\times150}=\frac{16}{65}
\]
\end{solution}

\begin{problem}
如图，在三棱锥 \(P - ABC\) 中，\(\triangle PAB\) 是等边三角形，\(\angle PAC = \angle PBC = 90^{\circ}\).

\begin{enumerate}
\item 求证：\(AB \perp PC\)；
\item 若 \(PC = 4\)，且平面 \(PAC \perp\) 平面 \(PBC\)，求三棱锥 \(P - ABC\) 体积.
\end{enumerate}

\bitmapfigure[max width=0.36\linewidth]{img/2009/new_curriculum_liberal/q18_fig1.png}
\end{problem}

\begin{solution}
设 \(P\) 为原点，记
\(\bs{a}=\overrightarrow{PA}\)，\(\bs{b}=\overrightarrow{PB}\)，\(\bs{c}=\overrightarrow{PC}\).
因为 \(\triangle PAB\) 是等边三角形，故 \(|\bs{a}|=|\bs{b}|\). 又由 \(\angle PAC=90^{\circ}\) 和 \(\angle PBC=90^{\circ}\)，有
\((-\bs{a})\cdot(\bs{c}-\bs{a})=0\)，\((-\bs{b})\cdot(\bs{c}-\bs{b})=0\)
即
\(\bs{a}\cdot\bs{c}=|\bs{a}|^2\)，\(\bs{b}\cdot\bs{c}=|\bs{b}|^2\).
而 \(\overrightarrow{AB}=\bs{b}-\bs{a}\)，所以
\[
\overrightarrow{AB}\cdot\overrightarrow{PC}=(\bs{b}-\bs{a})\cdot\bs{c}=|\bs{b}|^2-|\bs{a}|^2=0
\]
故 \(AB\perp PC\).

下求体积. 由两平面 \(PAC\perp PBC\)，可取 \(PC\) 为 \(z\) 轴，并使两个平面分别为 \(xz\) 平面和 \(yz\) 平面. 设 \(PC=4\)，并记 \(PA=PB=s\)，则可写成
\(A=(u,0,r)\)，\(B=(0,v,w)\)，\(C=(0,0,4)\).
由 \(PA\perp AC\) 得
\[
(-\overrightarrow{PA})\cdot\overrightarrow{AC}=(-u,0,-r)\cdot(-u,0,4-r)=0
\]
所以 \(u^2=r(4-r)\)，从而 \(s^2=u^2+r^2=4r\). 同理，由 \(PB\perp BC\) 得 \(s^2=4w\). 因此 \(r=w=s^2/4\).

等边三角形 \(PAB\) 的夹角 \(\angle APB=60^{\circ}\)，于是
\[
\overrightarrow{PA}\cdot\overrightarrow{PB}=|\overrightarrow{PA}|\,|\overrightarrow{PB}|\cos60^{\circ}=\frac{s^2}{2}
\]
另一方面 \(\overrightarrow{PA}\cdot\overrightarrow{PB}=rw\)，故
\[
\left(\frac{s^2}{4}\right)^2=\frac{s^2}{2}
\]
由于 \(s>0\)，解得 \(s^2=8\)，从而 \(r=w=2\)，且 \(u^2=v^2=8-2^2=4\). 适当选择 \(x\)，\(y\) 轴正方向，可取
\(A=(2,0,2)\)，\(B=(0,2,2)\)，\(C=(0,0,4)\).
所以
\[
V_{P-ABC}=\frac{1}{6}\left|\overrightarrow{PA}\cdot\left(\overrightarrow{PB}\times\overrightarrow{PC}\right)\right|=\frac{1}{6}\left|16\right|=\frac{8}{3}
\]
\end{solution}

\begin{problem}
某工厂有工人 \(1000\) 名，其中 \(250\) 名工人参加过短期培训（称为 A 类工人），另外 \(750\) 名工人参加过长期培训（称为 B 类工人）. 现用分层抽样方法（按 A 类，B 类分二层）从该工厂的工人中共抽查 \(100\) 名工人，调查他们的生产能力（此处生产能力指一天加工的零件数）.
\begin{enumerate}
\item 求甲，乙两工人都被抽到的概率，其中甲为 A 类工人，乙为 B 类工人；

\item 从 A 类工人中的抽查结果和从 B 类工人中的抽查结果分别如下表 1 和表 2.

\begin{center}
表 1：

\begin{tabular}{c|ccccc}
\hline
生产能力分组 & \([100,110)\) & \([110,120)\) & \([120,130)\) & \([130,140)\) & \([140,150)\) \\
\hline
人数 & \(4\) & \(8\) & \(x\) & \(5\) & \(3\) \\
\hline
\end{tabular}
\end{center}

\begin{center}
表 2：

\begin{tabular}{c|cccc}
\hline
生产能力分组 & \([110,120)\) & \([120,130)\) & \([130,140)\) & \([140,150)\) \\
\hline
人数 & \(6\) & \(y\) & \(36\) & \(18\) \\
\hline
\end{tabular}
\end{center}

\begin{enumerate}
\item 先确定 \(x\)，\(y\)，再完成下列频率分布直方图. 就生产能力而言，A 类工人中个体间的差异程度与 B 类工人中个体间的差异程度哪个更小？（不用计算，可通过观察直方图直接回答结论）
\item 分别估计 A 类工人和 B 类工人生产能力的平均数，并估计该工厂工人的生产能力的平均数. （同一组中的数据用该组区间的中点值作代表）
\end{enumerate}
\end{enumerate}

\texfigure[max width=0.90\linewidth]{img_repaint/2009/new_curriculum_liberal/q19_fig1.tex}
\FloatBarrier
\end{problem}

\begin{solution}
\begin{enumerate}
\item 分层抽样时，各层抽取人数与该层人数成比例. 因此 A 类、B 类分别抽取
\(100\times\frac{250}{1000}=25\)，\(100\times\frac{750}{1000}=75\)
人. 甲属于 A 类，乙属于 B 类，两层抽样相互独立，所以
\[
P=\frac{25}{250}\times\frac{75}{750}=\frac{1}{100}
\]

\item
\begin{enumerate}
\item 表 1 的总人数为 A 类抽样人数 \(25\)，表 2 的总人数为 B 类抽样人数 \(75\). 因而
\(4+8+x+5+3=25\)，\(6+y+36+18=75\)
解得
\(x=5\)，\(y=15\).

各组组距都为 \(10\)，所以每组柱高等于该组频率除以 \(10\). 表 1 五组柱高从左到右依次为
\(\frac{4/25}{10}=0.016\)，\(\frac{8/25}{10}=0.032\)，\(\frac{5/25}{10}=0.020\)，\(\frac{5/25}{10}=0.020\)，\(\frac{3/25}{10}=0.012\).
表 2 四组柱高从左到右依次为
\(\frac{6/75}{10}=0.008\)，\(\frac{15/75}{10}=0.020\)，\(\frac{36/75}{10}=0.048\)，\(\frac{18/75}{10}=0.024\).
据此在图中分别作出对应高度的矩形. 从两图中数据的分布情况可见，B 类工人的数据相对集中，故 B 类工人中个体间的差异程度更小.

\item A 类各组区间的中点分别为 \(105\)，\(115\)，\(125\)，\(135\)，\(145\)，所以
\[
\overline{x}_{A}=\frac{105\times4+115\times8+125\times5+135\times5+145\times3}{25}=\frac{3075}{25}=123
\]
B 类各组区间的中点分别为 \(115\)，\(125\)，\(135\)，\(145\)，所以
\[
\overline{x}_{B}=\frac{115\times6+125\times15+135\times36+145\times18}{75}=\frac{10035}{75}=133.8
\]
A 类、B 类工人分别占全厂工人的
\(\frac{250}{1000}=\frac14\)，\(\frac{750}{1000}=\frac34\).
因此全厂工人生产能力平均数的估计值为
\[
\overline{x}=\frac14\overline{x}_{A}+\frac34\overline{x}_{B}=\frac14\times123+\frac34\times133.8=131.1
\]
\end{enumerate}
\end{enumerate}
\end{solution}

\begin{problem}
已知椭圆 \(C\) 的中心为直角坐标系 \(xOy\) 的原点，焦点在 \(x\) 轴上，它的一个顶点到两个焦点的距离分别是 \(7\) 和 \(1\).

\begin{enumerate}
\item 求椭圆 \(C\) 的方程；
\item 若 \(P\) 为椭圆 \(C\) 上的动点，\(M\) 为过 \(P\) 且垂直于 \(x\) 轴的直线上的点，\(\frac{|OP|}{|OM|} = e\)（\(e\) 为椭圆 \(C\) 的离心率），求点 \(M\) 的轨迹方程，并说明轨迹是什么曲线.
\end{enumerate}
\end{problem}

\begin{solution}
\begin{enumerate}
\item 设椭圆的长半轴、半焦距分别为 \(a\)，\(c\). 由题意，某个顶点到两个焦点的距离为 \(a-c\) 和 \(a+c\)，所以
\(a-c=1\)，\(a+c=7\).
解得 \(a=4\)，\(c=3\). 因此短半轴的平方为
\[
b^2=a^2-c^2=16-9=7
\]
故椭圆 \(C\) 的方程为
\[
\frac{x^2}{16}+\frac{y^2}{7}=1
\]

\item 由 \(M\) 在过 \(P\) 且垂直于 \(x\) 轴的直线上，设 \(P=(x,y)\)，则 \(M=(x,m)\). 椭圆的离心率为
\[
e=\frac{c}{a}=\frac34
\]
题设条件给出
\(\frac{|OP|}{|OM|}=\frac34\)，\(x^2+y^2=\frac{9}{16}(x^2+m^2)\).
于是
\[
m^2=\frac79x^2+\frac{16}{9}y^2
\]
由椭圆方程 \(y^2=7-\dfrac{7}{16}x^2\)，代入上式得
\[
m^2=\frac79x^2+\frac{16}{9}\left(7-\frac{7}{16}x^2\right)=\frac{112}{9}
\]
所以 \(m=\pm\dfrac{4\sqrt7}{3}\). 又 \(P\) 在椭圆上时 \(-4\leqslant x\leqslant4\)，且每个这样的 \(x\) 都能取到椭圆上的 \(P\). 因此 \(M\) 的轨迹为
\(y=\frac{4\sqrt7}{3}\quad\text{或}\quad y=-\frac{4\sqrt7}{3}\)，\(-4\leqslant x\leqslant4\)
即两条平行于 \(x\) 轴的线段.
\end{enumerate}
\end{solution}

\begin{problem}
已知函数 \( f(x) = x^3 - 3ax^2 - 9a^2x + a^3 \).

\begin{enumerate}
\item 设 \(a = 1\)，求函数 \(f(x)\) 的极值；
\item 若 \(a > \frac{1}{4}\)，且当 \(x \in [1, 4a]\) 时，\(|f'(x)| \leqslant 12a\) 恒成立，试确定 \(a\) 的取值范围.
\end{enumerate}
\end{problem}

\begin{solution}
\begin{enumerate}
\item 当 \(a=1\) 时，
\(f(x)=x^3-3x^2-9x+1\)，\(f'(x)=3x^2-6x-9=3(x+1)(x-3)\).
当 \(x<-1\) 或 \(x>3\) 时 \(f'(x)>0\)，当 \(-1<x<3\) 时 \(f'(x)<0\). 因此 \(x=-1\) 为极大值点，\(x=3\) 为极小值点，并且
\(f(-1)=-1-3+9+1=6\)，\(f(3)=27-27-27+1=-26\).
所以函数的极大值为 \(6\)，极小值为 \(-26\).

\item 对一般的 \(a\)，有
\[
f'(x)=3x^2-6ax-9a^2=3(x-3a)(x+a)
\]
由 \(x\in[1,4a]\) 且 \(a>\dfrac14\)，可知区间非空. 先考察右端点 \(x=4a\)，有
\[
|f'(4a)|=|3(4a-3a)(4a+a)|=15a^2
\]
要使 \(|f'(x)|\leqslant12a\) 恒成立，必须有 \(15a^2\leqslant12a\). 因 \(a>0\)，得
\[
a\leqslant\frac45
\]

反之，设 \(\dfrac14<a\leqslant\dfrac45\)，则 \(a<1\)，从而在 \(x\in[1,4a]\) 上有 \(x\geqslant1>a\)，故 \(f'(x)\) 递增. 只需考察两个端点的绝对值. 右端点已经满足
\[
|f'(4a)|=15a^2\leqslant12a
\]
若 \(\dfrac14<a\leqslant\dfrac13\)，则 \(f'(1)=3(1-3a)(1+a)\geqslant0\)，且
\[
f'(1)\leqslant12a
\Longleftrightarrow 3a^2+6a-1\geqslant0
\]
由于 \(a>\dfrac14>\dfrac{2\sqrt3-3}{3}\)，上式成立. 若 \(\dfrac13\leqslant a\leqslant\dfrac45\)，则
\[
|f'(1)|=3(3a-1)(1+a)
\]
并且
\[
|f'(1)|\leqslant12a
\Longleftrightarrow 3a^2-2a-1\leqslant0
\Longleftrightarrow (3a+1)(a-1)\leqslant0
\]
这对 \(a\leqslant\dfrac45<1\) 成立. 因此该范围内确有 \(|f'(x)|\leqslant12a\) 恒成立.

综上，\(a\) 的取值范围为
\[
\left(\frac14,\frac45\right]
\]
\end{enumerate}
\end{solution}

\begin{problem}
如图，已知 \(\triangle ABC\) 的两条角平分线 \(AD\) 和 \(CE\) 相交于 \(H\)，\(\angle B = 60^{\circ}\)，\(F\) 在 \(AC\) 上，且 \(AE = AF\).

\begin{enumerate}
\item 证明：\(B\)，\(D\)，\(H\)，\(E\) 四点共圆；
\item 证明：\(CE\) 平分 \(\angle DEF\).
\end{enumerate}

\bitmapfigure[float=inline,max width=0.34\linewidth]{img/2009/new_curriculum_liberal/q22_fig1.png}
\FloatBarrier
\end{problem}

\begin{solution}
记 \(\angle A=A\)，\(\angle C=C\). 因为 \(\angle B=60^{\circ}\)，所以 \(A+C=120^{\circ}\).

\begin{enumerate}
\item \(D\) 在 \(BC\) 上，\(H\) 在 \(AD\) 上，故
\[
\angle BDH=\angle BDA=180^{\circ}-60^{\circ}-\frac{A}{2}=120^{\circ}-\frac{A}{2}
\]
又 \(E\) 在 \(AB\) 上，\(H\) 在 \(CE\) 上，且 \(CE\) 平分 \(\angle C\)，所以在 \(\triangle BEC\) 中
\[
\angle BEH=\angle BEC=180^{\circ}-60^{\circ}-\frac{C}{2}=120^{\circ}-\frac{C}{2}
\]
于是
\[
\angle BDH+\angle BEH=240^{\circ}-\frac{A+C}{2}=180^{\circ}
\]
因此四点 \(B\)，\(D\)，\(H\)，\(E\) 共圆.

\item \(AD\) 和 \(CE\) 是三角形的两条角平分线，故它们的交点 \(H\) 是内心，\(BH\) 也平分 \(\angle B\). 在圆 \(BDHE\) 中，\(\angle DEH\) 与 \(\angle DBH\) 同弧所对，所以
\[
\angle DEC=\angle DEH=\angle DBH=\frac12\angle B=30^{\circ}
\]
其中用到了 \(E\)，\(H\)，\(C\) 三点共线.

另一方面，\(AE=AF\)，所以 \(\triangle AEF\) 为等腰三角形. 由于 \(E\) 在 \(AB\) 上，\(F\) 在 \(AC\) 上，
\(\angle EAF=A\)，\(\angle AEF=\angle EFA=\frac{180^{\circ}-A}{2}=90^{\circ}-\frac{A}{2}\).
在 \(\triangle AEC\) 中，\(\angle EAC=A\)，\(\angle ACE=C/2\)，故
\[
\angle AEC=180^{\circ}-A-\frac{C}{2}
\]
由图中射线顺序 \(EA\)，\(EF\)，\(EC\) 可得
\[
\angle CEF=\angle AEC-\angle AEF=90^{\circ}-\frac{A+C}{2}=30^{\circ}
\]
于是 \(\angle DEC=\angle CEF\)，所以 \(CE\) 平分 \(\angle DEF\).
\end{enumerate}
\end{solution}

\FloatBarrier

\section{选考题：解答题}

\begin{problem}
已知曲线 \(C_1\)：\(\begin{cases}
x = -4 + \cos t,\\
y = 3 + \sin t,
\end{cases}\)（\(t\) 为参数），\(C_2\)：\(\begin{cases}
x = 8\cos \theta,\\
y = 3\sin \theta.
\end{cases}\)（\(\theta\) 为参数）.

\begin{enumerate}
\item 化 \(C_1\)，\(C_2\) 的方程为普通方程，并说明它们分别表示什么曲线；
\item 若 \(C_1\) 上的点 \(P\) 对应的参数为 \(t = \frac{\pi}{2}\)，\(Q\) 为 \(C_2\) 上的动点，求 \(PQ\) 中点 \(M\) 到直线 \(C_3\)：\(\begin{cases}
x = 3 + 2t,\\
y = -2 + t,
\end{cases}\)（\(t\) 为参数）距离的最小值.
\end{enumerate}
\end{problem}

\begin{solution}
\begin{enumerate}
\item 对 \(C_1\)，由 \(x+4=\cos t\)，\(y-3=\sin t\)，平方相加得
\[
(x+4)^2+(y-3)^2=\cos^2t+\sin^2t=1
\]
所以 \(C_1\) 是圆心为 \((-4,3)\)、半径为 \(1\) 的圆.

对 \(C_2\)，由 \(x=8\cos\theta\)，\(y=3\sin\theta\)，得
\[
\frac{x^2}{64}+\frac{y^2}{9}=\cos^2\theta+\sin^2\theta=1
\]
所以 \(C_2\) 是以原点为中心，长半轴为 \(8\)、短半轴为 \(3\) 的椭圆.

\item 当 \(t=\dfrac\pi2\) 时，\(P=(-4,4)\). 设 \(Q=(8\cos\theta,3\sin\theta)\)，则线段 \(PQ\) 的中点为
\[
M=\left(\frac{8\cos\theta-4}{2},\frac{3\sin\theta+4}{2}\right)=\left(4\cos\theta-2,\frac32\sin\theta+2\right)
\]
直线 \(C_3\) 的参数方程为 \(x=3+2t\)，\(y=-2+t\). 消去参数 \(t\)，得
\[
x-2y-7=0
\]
点 \(M\) 到该直线的距离为
\[
d=\frac{|x_M-2y_M-7|}{\sqrt5}=\frac{|4\cos\theta-3\sin\theta-13|}{\sqrt5}
\]
由 \(4\cos\theta-3\sin\theta\leqslant\sqrt{4^2+(-3)^2}=5\)，且该最大值可在 \((\cos\theta,\sin\theta)=\left(\dfrac45,-\dfrac35\right)\) 时取得，故括号内始终为负，且
\[
d\geqslant\frac{13-5}{\sqrt5}=\frac{8}{\sqrt5}=\frac{8\sqrt5}{5}
\]
当 \((\cos\theta,\sin\theta)=\left(\dfrac45,-\dfrac35\right)\) 时取等号，所以最小距离为 \(\dfrac{8\sqrt5}{5}\).
\end{enumerate}
\end{solution}

\begin{problem}
如图，\(O\) 为数轴的原点，\(A\)，\(B\)，\(M\) 为数轴上三点，\(C\) 为线段 \(OM\) 上的动点，设 \(x\) 表示 \(C\) 与原点的距离，\(y\) 表示 \(C\) 到 \(A\) 距离 \(4\) 倍与 \(C\) 到 \(B\) 距离的 \(6\) 倍的和.

\begin{enumerate}
\item 将 \(y\) 表示成 \(x\) 的函数；
\item 要使 \(y\) 的值不超过 \(70\)，\(x\) 应该在什么范围内取值.
\end{enumerate}

\FigureLayoutDeclare{img/2009/new_curriculum_liberal/q24_fig1.png}{8.0cm}{23bp 58bp 77bp 37bp}
\bitmapfigure[max width=0.46\linewidth]{img/2009/new_curriculum_liberal/q24_fig1.png}
\end{problem}

\begin{solution}
由图知 \(O\)，\(A\)，\(B\)，\(M\) 的坐标依次为 \(0\)，\(10\)，\(20\)，\(30\). 由于 \(C\) 在线段 \(OM\) 上，故
\[
0\leqslant x\leqslant30
\]
点 \(C\) 到 \(A\)，\(B\) 的距离分别为 \(|x-10|\)，\(|x-20|\)，所以
\[
y=4|x-10|+6|x-20|
\]
按 \(x\) 与 \(10\)，\(20\) 的大小关系分段，得到
\[
y=\begin{cases}
160-10x,&0\leqslant x<10,\\
80-2x,&10\leqslant x<20,\\
10x-160,&20\leqslant x\leqslant30.
\end{cases}
\]
现求 \(y\leqslant70\). 在第一段，\(160-10x\leqslant70\) 等价于 \(x\geqslant9\)，得到 \(9\leqslant x<10\)；在第二段，\(80-2x\leqslant70\) 等价于 \(x\geqslant5\)，故整段 \(10\leqslant x<20\) 均满足；在第三段，\(10x-160\leqslant70\) 等价于 \(x\leqslant23\)，得到 \(20\leqslant x\leqslant23\). 合并三段并补上分界点，得
\[
x\in[9,23]
\]
\end{solution}
